\documentclass[12pt]{article}
\usepackage[T1]{fontenc}
\usepackage[utf8]{inputenc}
\usepackage{lmodern}
\usepackage{textcomp}
\usepackage{enumitem}
\usepackage{geometry}
\geometry{margin=1in}
\begin{document}
\begin{center}
Answer Key - Biology - Undergraduate Level Assessment 17 \
\small (Answers are ordered exactly as in the question paper.)
\end{center}

\section*{Multiple Choice}
\begin{enumerate}[label=\arabic*.]
\item \textbf{Answer:} A
\\ \textit{Options:} A) 108 nucleotides; B) 120 nucleotides; C) 150 nucleotides; D) 144 nucleotides; E) 60 nucleotides
\\ \textit{Note:} The correct answer is 108 nucleotides because in a DNA molecule, adenine (A) pairs with thymine (T) and cytosine (C) pairs with guanine (G), so if 20\% of the 180 base pairs are adenine, there are 36 adenine nucleotides, which means there are also 36 thymine nucleotides, leaving 108 nucleotides for cytosine and guanine combined, and since they must be equal, there are 54 cytosine nucleotides.
\item \textbf{Answer:} A
\\ \textit{Options:} A) Cortisol, aldosterone, and adrenal androgens; B) Insulin, glucagon, and somatostatin; C) Cholecystokinin, ghrelin, and leptin; D) Vasopressin, erythropoietin, and renin; E) Epinephrine, norepinephrine, and dopamine
\\ \textit{Note:} Cortisol, aldosterone, and adrenal androgens are key hormones produced by the adrenal glands that play crucial roles in the regulation of growth, metabolism, and the development of secondary sexual characteristics during puberty in both males and females.
\item \textbf{Answer:} E
\\ \textit{Options:} A) Epilepsy is the result of a chromosomal aneuploidy, similar to Down syndrome.; B) Epilepsy is always inherited in an X-linked recessive pattern.; C) Epilepsy is a dominant trait.; D) Epilepsy is caused by an autosomal dominant trait.; E) Epilepsy is due to an autosomal recessive trait.
\\ \textit{Note:} The correct answer is based on the understanding that certain forms of epilepsy can be inherited through an autosomal recessive pattern, where the disease manifests when an individual inherits two copies of the mutated gene, highlighting the genetic basis of the condition.
\item \textbf{Answer:} E
\\ \textit{Options:} A) Lysosome; B) Chloroplast; C) Ribosome; D) Golgi apparatus; E) Plasma membrane
\\ \textit{Note:} The plasma membrane is considered to have arisen first in the formation of the earliest cells because it provides a necessary boundary that allows for the compartmentalization of biochemical processes, thereby facilitating the emergence of life by enabling the maintenance of an internal environment distinct from the external surroundings.
\item \textbf{Answer:} D
\\ \textit{Options:} A) Centrioles; B) Chloroplasts; C) Lysosomes; D) Cell wall; E) Mitochondria
\\ \textit{Note:} The presence of a cell wall in prokaryotic cells, which is typically composed of peptidoglycan, distinguishes them from animal cells that lack a cell wall and instead have only a plasma membrane.
\end{enumerate}

\section*{True / False}
\begin{enumerate}[label=\arabic*.]
\setcounter{enumi}{5}
\item \textbf{Answer:} True
\\ \textit{Options:} A) True; B) False
\\ \textit{Note:} Polygenic traits are true because they arise from the cumulative effects of multiple genes, each contributing to the overall phenotype, resulting in a continuous range of variations rather than discrete categories.
\item \textbf{Answer:} True
\\ \textit{Options:} A) True; B) False
\\ \textit{Note:} Cycads are unique among seed plants because they have a reproductive strategy that allows their seeds to remain dormant for extended periods before germination, adapting to environmental conditions.
\item \textbf{Answer:} True
\\ \textit{Options:} A) True; B) False
\\ \textit{Note:} Complete equilibrium in a gene pool is not expected because factors such as mutation introduce new genetic variations, while natural selection favors certain traits over others, continually altering allele frequencies within a population.
\item \textbf{Answer:} True
\\ \textit{Options:} A) True; B) False
\\ \textit{Note:} In chick embryos, the yolk sac is crucial for nourishment as it contains the yolk, which provides essential nutrients for the developing embryo during the early stages of development before it can obtain food from the environment.
\item \textbf{Answer:} False
\\ \textit{Options:} A) True; B) False
\\ \textit{Note:} Distantly related organisms can exhibit significant physical differences despite sharing a common ancestor due to divergent evolution, environmental adaptations, and variations in selective pressures over time.
\end{enumerate}

\section*{Long Form Response}
\begin{enumerate}[label=\arabic*.]
\setcounter{enumi}{10}
\item \textbf{Answer:} In an animal with a diploid number of 8, the number of chromatids varies at different stages of meiosis. (a) During the tetrad stage of meiosis, each homologous chromosome pairs with its counterpart, resulting in a total of 4 chromatids, as each chromosome consists of two sister chromatids. (b) By late telophase of the first meiotic division, the cell has separated its homologous chromosomes, yielding 12 chromatids in total, since each of the original 4 chromosomes has been duplicated during DNA replication prior to meiosis. (c) In metaphase of the second meiotic division, the chromatids align at the metaphase plate, and there are 6 chromatids present, as the two cells produced from the first meiotic division each contain 3 chromosomes, each still composed of two chromatids.
\\ \textit{Note:} correct because it accurately describes the number of chromatids at each specified stage of meiosis, reflecting the processes of chromosome pairing, separation, and chromatid duplication in an organism with a diploid number of 8.
\item \textbf{Answer:} The characteristics of organisms within an ecosystem are significantly influenced by several environmental factors, including temperature, precipitation, soil chemistry, and the presence of decomposers. Temperature affects metabolic rates and growth patterns; for example, warmer climates often support a greater diversity of species, as seen in tropical rainforests compared to temperate forests. Precipitation determines water availability, influencing plant types and, consequently, herbivore populations-deserts illustrate this with their specialized drought-resistant flora. Soil chemistry, including pH and nutrient content, shapes plant communities, which in turn affects the entire food web; for instance, the rich, nutrient-dense soils of river deltas foster lush vegetation and diverse animal life. Lastly, the types of decomposers present play a crucial role in nutrient cycling, with specific decomposers like fungi and bacteria breaking down organic matter, thereby supporting plant growth and overall ecosystem health.
\\ \textit{Note:} correct because it identifies and explains the impact of critical environmental factors, such as temperature and precipitation, on organism characteristics and provides relevant examples that illustrate these influences within different ecosystems.
\end{enumerate}

\vfill
\noindent\textit{End of Answer Key}
\end{document}